A \emph{retraction} of a space $X$ onto a subset $A\subset X$ is a map
$r:X\rightarrow A$ such that $r|_A=\mathrm{id}$. In each case decide
whether or not there is a retraction of the pictured surface $X$ onto
the given subset $A$. Describe the retraction if it exists. Otherwise,
prove that there is no retraction. 
%%%%%%%%%%%%%%%%%%%%%%%%%%%%%%%%%%%%%%%%%%%%%%%%%%%%%%%%%%%%%%%%%%%
\begin{alist}
%------------------------------------------------------------------*
\item $A$ is a circle that wraps around the middle of a two-holed 
torus $X$, and separates the two-holed torus into two once-punctured
tori.
%------------------------------------------------------------------*
\item $A$ is the typical wedge of two circles seen as a subset of a
torus $X$.
%------------------------------------------------------------------*
\item $A$ is made up of two circles that separate a two-holed torus into
a four-times-punctured sphere.
%------------------------------------------------------------------*
\end{alist}
%%%%%%%%%%%%%%%%%%%%%%%%%%%%%%%%%%%%%%%%%%%%%%%%%%%%%%%%%%%%%%%%%%%
